\documentclass[12pt]{article}
\usepackage{amsmath,amssymb}

\begin{document}
\section*{{2022 AIME I Problems/Problem 15}}
Source: AoPS Wiki

\subsection*{Solution}
First, let define a triangle with side lengths $\sqrt{2x}$ , $\sqrt{2z}$ , and $l$ , with altitude from $l$ 's equal to $\sqrt{xz}$ . $l = \sqrt{2x - xz} + \sqrt{2z - xz}$ , the left side of one equation in the problem. Let $\theta$ be angle opposite the side with length $\sqrt{2x}$ . Then the altitude has length $\sqrt{2z} \cdot \sin(\theta) = \sqrt{xz}$ and thus $\sin(\theta) = \sqrt{\frac{x}{2}}$ , so $x=2\sin^2(\theta)$ and the side length $\sqrt{2x}$ is equal to $2\sin(\theta)$ . We can symmetrically apply this to the two other equations/triangles. By law of sines, we have $\frac{2\sin(\theta)}{\sin(\theta)} = 2R$ , with $R=1$ as the circumradius, same for all 3 triangles.
The circumcircle's central angle to a side is $2 \arcsin(l/2)$ , so the 3 triangles' $l=1, \sqrt{2}, \sqrt{3}$ , have angles $120^{\circ}, 90^{\circ}, 60^{\circ}$ , respectively. This means that by half angle arcs, we see that we have in some order, $x=2\sin^2(\alpha)$ , $y=2\sin^2(\beta)$ , and $z=2\sin^2(\gamma)$ (not necessarily this order, but here it does not matter due to symmetry), satisfying that $\alpha+\beta=180^{\circ}-\frac{120^{\circ}}{2}$ , $\beta+\gamma=180^{\circ}-\frac{90^{\circ}}{2}$ , and $\gamma+\alpha=180^{\circ}-\frac{60^{\circ}}{2}$ . Solving, we get $\alpha=\frac{135^{\circ}}{2}$ , $\beta=\frac{105^{\circ}}{2}$ , and $\gamma=\frac{165^{\circ}}{2}$ . We notice that \[[(1-x)(1-y)(1-z)]^2=[\sin(2\alpha)\sin(2\beta)\sin(2\gamma)]^2=[\sin(135^{\circ})\sin(105^{\circ})\sin(165^{\circ})]^2\] \[=\left(\frac{\sqrt{2}}{2} \cdot \frac{\sqrt{6}-\sqrt{2}}{4} \cdot \frac{\sqrt{6}+\sqrt{2}}{4}\right)^2 = \left(\frac{\sqrt{2}}{8}\right)^2=\frac{1}{32} \to \boxed{033}. \blacksquare\] - kevinmathz

\end{document}
